% ======================================================================
% 1. INTRODUCTION  (budget: 1.5 pages)
% Goal: the reader accepts that "password" is the wrong unit of analysis
% before chapter 2 starts defining terms.
% ======================================================================
\section{Introduction}
\label{sec:intro}

\subsection{Motivation}
% TODO
% - Authentication protocols are usually attacked at their design assumptions,
%   not their ciphers. NTLM is a case study: no cryptographic break exists,
%   and yet the protocol is trivially defeated when the wrong secret leaks.
% - Scale of exposure: NTLM is still negotiable across SMB, RPC, WinRM, LDAP,
%   MSSQL and HTTP/Negotiate; AD domain controllers still accept RC4 tickets.
% - Microsoft has announced NTLM's deprecation -- use this as the framing hook:
%   the technique matters most precisely while the protocol is still in place.

\subsection{Problem statement}
% TODO: the single sentence the essay defends. Suggested:
% "Because both protocols allow a client to complete authentication using a
%  locally stored derived key, any disclosure of that key is equivalent to
%  disclosure of the account, regardless of password strength."

\subsection{Contributions}
% TODO: numbered list, 4 items
%   1. A from-first-principles account of the key material Windows holds.
%   2. A derivation-level explanation of why PtH is not replay.
%   3. A reproducible PoC (PtH over SMB/WinRM, PtK to TGT) with captured evidence.
%   4. Measured before/after of the defences that actually stop it.
% Link the repository here (also in the title block).

\subsection{Scope and non-goals}
% TODO: one short paragraph. State that ticket and certificate abuse share the
% same root cause but are out of scope, and say why (different lab requirements).
% Keep this brief: one or two sentences, no defensive justification.
